\documentclass[]{article}

\usepackage{amsmath}
\usepackage{multirow}
\usepackage{natbib}
\usepackage{graphicx} 


\begin{document}

% Auto-generated figures file
% \usepackage{graphicx}

\begin{figure}[h!]
    \centering
    \includegraphics[width=0.5\textwidth]{../input_files/plots/step_5_nll_histogram_1_20260406_123725.png}
    \caption{Distribution of the Negative Log-Likelihood (NLL) Out-of-Distribution (OoD) scores evaluated on the 10,000 test maps. The left panel shows that the bulk of the samples are concentrated at low NLL values, consistent with the in-distribution data. The right panel displays the full distribution on a log-log scale, revealing a heavy tail of samples with extremely high NLL scores. This long tail identifies the anomalous OoD maps, whose non-Gaussian small-scale structures are assigned very low probability by the conditional density model.}
    \label{fig:image0}
\end{figure}

\begin{figure}[h!]
    \centering
    \includegraphics[width=0.5\textwidth]{../input_files/plots/step_6_roc_curve_20260406_124510.png}
    \caption{Validation Receiver Operating Characteristic (ROC) curve for the synthetic Out-of-Distribution (OoD) detection task, where OoD samples were generated by rescaling the signal in convergence maps. The curve shows the True Positive Rate (TPR) versus the False Positive Rate (FPR) on a logarithmic scale. The shaded region highlights the critical FPR range of [0.001, 0.05] over which the performance metric is evaluated. The steep rise of the curve within this region indicates that the model effectively identifies OoD samples while maintaining a false alarm rate below 5\%, demonstrating high sensitivity to structural deviations.}
    \label{fig:image1}
\end{figure}



\end{document}
                